
\documentclass[a4paper]{article}

\usepackage[utf8]{inputenc}   % LaTeX, comprends les accents !
\usepackage[T1]{fontenc}      % Police contenant les caractères français
\usepackage{geometry}         % Définir les marges
\usepackage[francais]{babel}  % Placez ici une liste de langues, la
                              % dernière étant la langue principale
\usepackage{graphicx}         % Inserer des images
\usepackage{gensymb}
\usepackage{color}

\pagestyle{headings}          % Pour mettre des entêtes avec les titres
                              % des sections en haut de page

\title{\bf The NIKA polarization pipeline}           % Les paramètres du titre : titre, auteur, date
\author{{\it The NIKA collaboration}}
\date{\today}

\setlength{\textwidth}{15.5cm}
\setlength{\textheight}{23cm}
\oddsidemargin +0.5cm                                          % Marges
\evensidemargin +0.5cm
\topmargin +0cm

\begin{document}

\maketitle                    % Faire un titre utilisant les données
                              % passées à \title, \author et \date
\begin{figure}[!h]
\centering
\includegraphics[height=2cm,trim=0cm 0.1cm 0.1cm 0.1cm, clip=true]{Figures/nika_white_bg}
\end{figure}
\vspace*{0.5cm}

\setcounter{page}{1}
\setcounter{tocdepth}{1}

\section{Work plan}

\subsection{Nov. 18th}

\begin{itemize}
\item Alessia is looking at the lab data and trying to get consistent
  estimates of the tranmission coefficients of the HWP $\alpha$ and $\beta$ and
  the phase shift angle $\phi$. She will then put these parameters in the
  simulation pipeline to determine the impact of these nuisance parameters if they
  are ignored in the map-making. She will then see how to correct for them and set
  constraints on the precision we need on their determination.
\item Alessia will vary the scan speed, the HWP rotation frequency and the
  sampling rate to determine the possible working points of the system that
  minimize the template and the duration of observations
\item These two items will lead to the definition of calibration actions we
  will have to perform at Pico in January.
\item Nico will continue analyzing the data taken during the Run Cryo and
  maintain compatibility of the codes.
\end{itemize}



\subsection{Scan speed / HWP rotation frequency / Acquisision rate: TBD}

The signal is contained in the frequency range allowed by the optical beam and
the scanning frequency. We want it to be small compared to the HWP rotation
frequency because... On the other hand, we do not want to scan too slowly to
avoid too long integration time and the HWP rotation frequency impacts the
acquisition rate.\\

We need to determine the specifications to put on each of these parameters.


\subsection{Instrumental polarization: TBD}

Definition: Instrumental Polarization (IP) is the ability of the experimental setup to induce
polarization from an unpolarized incoming signal before the modulation by HWP.\\

How much can we tolerate ? How much can we correct for and how ? Propagation of
errors with a fiducial Mueller matrix should give first estaimtes.\\

We should distinguish between IP acting on the near field (namely the
atmosphere) and IP acting on astronimcal signal. We may address both separately
with skydips and OTF geometries on a planet at different elevations.\\

What about a polarized background coming from the Nasmyth cabin and entering
directly in the cryostat ?

\subsection{HWP acquisition: TBD}

The HWP rotation angle is not sampled at the same rate as the data rate. How is
it done exactly ? How does it impact the measured signal ?

\subsection{HWP modeling: TBD}

For now, the simulation and analysis codes assume a perfect HWP. We should make
this more physical.

\subsection{Decorrelation / Atmosphere: TBD}

Nov.~12th, 2013: the level of sky noise is arbitrary for now and has the same
amplitude at 1 and 2mm. We should look at how it's actually calibrated in the
temperature pipeline and reuse it.\\

How does the atmosphere impact on the recovery of $S_0, S_1, S_2$ ? How do we
decorrelate ?

\subsection{Transfer function: TBD}

Any processing applied to the data alters the incoming signal. In particular,
the current polarizatin map making includes a linear template subtraction and a
lowpass. How does it alter the signal ? (run simulations of various input skies,
in particular power laws with known spectra).

\subsection{Intesnsity and polarizaiton bands: TBD}

The bands are not the same in intensity and polarization. for now, the
simulation and analysis pipelines ignore this issue.


\subsection{Template: TBD}

Def: we call template any additional signal to the timelines that is due to the
HWP. It should be a function of the HWP angle and should look at first order
like a sum or harmonics of the HWP rotation frequency with slowly varying
amplitudes.\\

A piece of code can fit for this template and subtract it. For now it deals with
one detector at a time. It could be improved by using several detectors at the
same time (TBC).\\

Can we go towards a model of this template with more physics in it to subtract
it better ?


\subsection{Astrometry: TBD}

\subsubsection{Parallactic angle: TBD}
Nov. 12th, 2013: As the elevation of the source changes, its polarization rotates too. For now,
the simulation and the analysis codes do not account for this.

\subsubsection{Calibration of the orientation:TBD}
Nov. 12th, 2013: For now, we implicitely assme that the zero of the HWP angle coincides with the
zero on the sky (sic!). How will we actually derive the phase of the HWP angle ?

\subsection{Beams: TBD}
Nov. 12th, 2013: For now, the simulation pipeline assumes that all kids have the same beam. This
is to avoid to produce one map per different kid beam. A FICSbell approach could
help improve the simulation.

\subsection{Splitting grid: TBD}
Nov. 12th, 2013: The splitting grid is assumed to be perfect for now.

\subsection{Pixel aliasing: TBD}
Nov. 12th, 2013: The diffuse component of the simulated sky is obtained by
scanning a simulated pixelized map. If the pixels are not small enough compared
to the sampling rate and scan speed, this induces artefacts in the timelines
that hamper our ability to compare input/output. For now, the only solution is
to simulate input maps with much small pixels but then comparison with output
map with coarser resolution is not trivial.

\end{document}

